\documentclass[11pt,letterpaper]{article}
\usepackage[T1]{fontenc}
\usepackage[utf8]{inputenc}
\usepackage{lmodern,microtype}
\usepackage[margin=0.9in,headheight=15pt]{geometry}
\usepackage{amsmath,amssymb,amsthm,mathtools}
\usepackage{booktabs,longtable,array,enumitem}
\usepackage{xcolor,xurl,hyperref,fancyhdr}
\definecolor{ink}{HTML}{19334C}
\hypersetup{colorlinks=true,linkcolor=ink,citecolor=ink,urlcolor=ink,
 pdftitle={Single-Exponential Activation Kernels for Dual-Forest Selection},
 pdfauthor={Research contribution prepared with ChatGPT for the ProveIt project}}
\pagestyle{fancy}\fancyhf{}
\fancyhead[L]{\small\sffamily Single-exponential activation kernels}
\fancyhead[R]{\small\sffamily ProveIt research continuation}
\fancyfoot[C]{\thepage}
\setlength{\parskip}{4pt}\setlength{\parindent}{1.15em}
\setlength{\emergencystretch}{2em}
\setlist[enumerate]{itemsep=3pt,topsep=4pt}
\newtheorem{theorem}{Theorem}[section]
\newtheorem{lemma}[theorem]{Lemma}
\newtheorem{proposition}[theorem]{Proposition}
\newtheorem{corollary}[theorem]{Corollary}
\theoremstyle{definition}
\newtheorem{definition}[theorem]{Definition}
\newtheorem{example}[theorem]{Example}
\newtheorem{remark}[theorem]{Remark}
\newcommand{\F}{\mathbb F}
\newcommand{\Reach}{\operatorname{Reach}}
\newcommand{\cl}{\operatorname{cl}}
\newcommand{\supp}{\operatorname{supp}}
\newcommand{\rk}{\operatorname{rank}}
\newcommand{\poly}{\operatorname{poly}}
\newcommand{\Bell}{\operatorname{Bell}}
\newcommand{\pin}{\texttt{e0b2f802f336fcb51e9ec5f29174cfe28dfb9d39}}
\newcommand{\Yes}{\texttt{SUCCESS}}
\newcommand{\No}{\texttt{BLOCKED\_EXACT}}
\newcommand{\Unk}{\texttt{UNKNOWN}}
\begin{document}
\hypersetup{pageanchor=false}
\begin{titlepage}
\vspace*{0.15in}
{\large\sffamily\color{ink} TOPOLOGY / UNKNOT RECOGNITION\par}
\vspace{0.35in}
{\Huge\bfseries\color{ink} Single-Exponential\par Activation Kernels\par}
\vspace{0.15in}
{\LARGE\bfseries for Dual-Forest Selection\par}
\vspace{0.27in}
{\large Generic mandatory contraction, representative families,\par
and deterministically checkable low-rank obstructions\par}
\vspace{0.3in}
{Research contribution prepared with ChatGPT\par
for Vladimir Reshetnikov's \emph{ProveIt} project\par}
{12 October 2026\par}
\vspace{0.25in}
\begin{abstract}
The audited incoming ProveIt dual-forest selector has a deterministic
$2^{O(b\log(b+1))}\poly(\Lambda)$ bound, where $b$ counts rows incident to
usable higher-support columns and $\Lambda$ is the encoded local input size.
We give a one-sided randomized $2^{O(b)}\poly(\Lambda,\lambda)$ search
with miss probability at most $2^{-\lambda}$ on positive inputs. Every
reported success is an independently replayable, entirely discrete certificate.
The proof replaces exceptional-row permutations and shared-witness partitions
by a closed activation state, at most $b$ useful additional seeds, generic
contraction of mandatory color representatives, and exterior-algebra
representative families. Only the dual graph is represented over a finite
field; geometric intersection counts are never reduced modulo that field.

When there are no mandatory low-column source components and the dual graphic
rank is at most $b$, the same method is deterministic and exact. A negative
answer has a frontier-span certificate checked from source graph minors and
literal peeling, without trusting the optimizer. We also give a necessary
source-chain parity condition and a polynomial criterion for parity-compatible
lifts of saturated intersection data.

The package contains 71 passing unit tests, 8,061 exhaustive/random
matrix--graph comparisons, 300 planted positive trials, independent
finite-field audits, comparative timings, and an opt-in record-level native
adapter. All 3,417 positive trials were certified; 3,391 negative inputs
had exact local certificates and 1,553 remained explicitly inconclusive.
These results improve a specified local selection stage. They do not establish
a general quasi-polynomial unknot recognizer or a geometric realization theorem
for the abstract benchmark families.
\end{abstract}
\vfill
{\small\textbf{Status.} Conventional proofs and executed research code; not
externally refereed or proof-assistant formalized. The contribution is measured
against the audited predecessor, not claimed as a historically first use of
representative families. Two other recently incoming archives were identified
but were not content-audited.}\par
\end{titlepage}
\hypersetup{pageanchor=true}\pagenumbering{roman}
\tableofcontents
\clearpage\pagenumbering{arabic}

\section{Research position and the advance}
\subsection{The checkpoint being continued}
This study fixes the maintained repository at commit
\begin{center}\small\pin.\end{center}
The source handles, dual endpoints, preservation of bundle ports, and the
five-row loop experiment are taken from the inspected maintained materials
\cite{repo,dual,loop,kernel}. The immediate algorithmic predecessor is
\emph{Matroidal Dual-Forest Selection} \cite{prior}, supplied in the incoming
archive of 12 October 2026. Its higher-support algorithm considers exceptional
row orders and partitions of the demands that may share a seed. The bound is
$b!\Bell(b)\poly(\Lambda)$, hence $2^{O(b\log(b+1))}\poly(\Lambda)$.
Its question on replacing this dependence by $2^{O(b)}$ is the concrete target
addressed here.

The distinction between choosing \emph{some} boundary-rooted dual forest and
choosing a \emph{useful} one is genuine. In the maintained loop example, the
default tree deletes a row that leaves a stopping set, while three alternative
choices allow completion \cite{loop}. The prior selector already resolves
that coverage issue. Our reconstructed loop experiment is consequently a
performance and certificate-replay comparison, not a claim to have newly
covered that source.

The ingredients from linear matroid representative-family theory are
classical; useful algorithmic accounts include \cite{fomin,bevern}. Our
contribution is the precise reduction of this activation problem to a small
representative-state search, including how mandatory source components are
handled without retaining all their choices. Generic linear combinations are
used only as a search device. The final tree and peeling schedule do not
contain field elements or symbolic variables.

\subsection{What changes, and what does not}
Write $r=|V(D)|-1$ for the rank of the connected dual graph, and let $q$ be the
number of mandatory components defined in Section~\ref{sec:activation}. After
an optional exact mandatory-obstruction prepass, put
\[
 d=\min\{b,r-q\}.
\]
One random round has the conservative field-arithmetic bound
\begin{equation}\label{eq:headline}
 \poly(\Lambda)+2^b8^d\poly(\Lambda)
 \ \leq\ 16^b\poly(\Lambda),
\end{equation}
with single-exponential bit complexity as well. Repeating $\lambda$ times
makes the miss probability at most $2^{-\lambda}$ on any input that has a
successful tree. The dependence on the number $q$ of mandatory components is
polynomial, not exponential. When $q=0$ and $r\leq b$, the original binary
incidence representation is already small enough: the search is deterministic
and can certify either outcome.

The input remains a fixed unsigned core matrix and a fixed dual graph, with
the original zero/one-handle forest frozen. The method does not optimize that
primal forest, discover a normal surface, establish admissibility, or maintain
all boundary-pattern data through an entire hierarchy. A negative local
certificate says that \emph{this particular tree-plus-positive-cancellation
strategy} cannot succeed. It says nothing by itself about whether the source
is a handlebody or whether a knot is trivial.

Lackenby's announced target is $n^{O(\log n)}$ in the crossing number
\cite{announcement}. The inspected 2026 hierarchy paper distinguishes its
iteration accounting from the still unspecified costs of some operations in
Section~9 \cite{lackenby}. The local quasi-polynomial regime proved below is
therefore stated with explicit source-size, call-count, and feasibility
hypotheses, rather than presented as a complete recognition bound.

\section{The fixed cancellation problem}\label{sec:model}
\subsection{Unsigned counts and dual edges}
Let $E=\{0,\ldots,m-1\}$ be a row set, let $C$ be a column set, and let
\[
 N=(n_{eh})_{e\in E,h\in C}\in\mathbb Z_{\geq0}^{E\times C}.
\]
Rows model source two-handles, and columns model the one-handles remaining
after the chosen zero/one forest contraction. Counts are unsigned geometric
intersections when a geometric source is supplied. For the abstract selection
problem only the categories $0$, $1$, and at least $2$ matter.

Let $D=(V,E)$ be a connected multigraph with a distinguished boundary root
$0$. The same label $e$ denotes a matrix row and a dual edge. Parallel edges
and loops are retained; a loop is not a forest edge. Let $r=|V|-1$. A spanning
tree $T$ first deletes its row set from $N$. Thereafter a row $e$ can be peeled
using column $h$ exactly when
\begin{equation}\label{eq:private}
 n_{eh}=1,\qquad n_{fh}=0\quad\text{for every other live row }f.
\end{equation}
Both the row and that pivot column are then removed. A \emph{successful tree}
is a spanning tree after whose deletion all remaining rows can be peeled.

The encoded input size $\Lambda$ includes the graph, column information,
matrix entries, and the bit lengths of the supplied counts. All graph and
matrix objects used here are explicit. The implementation retains original
counts for input binding, while tests of $n=0$ and $n=1$ implicitly saturate
larger counts. Nothing is proportional to the numerical value of a large
intersection count.

\subsection{Why static peeling is the correct local target}
The maintained cancellation theorem and the audited predecessor show that
private-column peeling is complete for the positive unsigned update model
\cite{peeling,prior}. We recall the stopping-set argument so that the local
criterion is unambiguous.

\begin{definition}
A nonempty set of rows $R$ is a \emph{stopping set} if, for every column $h$,
\[
 S_h(R):=\sum_{e\in R}n_{eh}\quad\text{is either }0\text{ or at least }2.
\]
\end{definition}
A fixed-row deletion set succeeds under peeling if and only if it meets every
stopping set of the original matrix. Indeed, a peeling process cannot remove
the first row of a surviving stopping set. Conversely, the nonempty terminal
set of any stalled peeling process is a stopping set. Thus the terminal set
is independent of the legal peeling order: it is the union of all stopping
sets avoiding the initial deletions.

For completeness with respect to the positive-fill model, an unsaturated
unit pivot at $(e,g)$ performs
\[
 n'_{fh}=n_{fh}+n_{fg}n_{eh}\qquad(f\ne e,\ h\ne g).
\]
If $e\in R$ and $n_{eg}=1$, then
\begin{equation}\label{eq:stopping-update}
 S'_h(R\setminus\{e\})
   =S_h(R)+(S_g(R)-2)n_{eh}.
\end{equation}
Since $S_g(R)\geq2$, no surviving stopping-set column acquires total one.
The remaining row set is nonempty because $e$ contributes only one to column
$g$. A pivot outside $R$ cannot use a column of total zero on $R$ to create
an entry there, and otherwise only increases its entries. Stopping sets
therefore persist through arbitrary positive pivots. Saturating each entry
at two preserves precisely the distinction between a column total of zero,
one, and at least two. This proves that a stalled private-column process
cannot be repaired merely by choosing other positive-fill pivots.

This argument does \emph{not} apply to arbitrary signed changes of attaching
words, geometric slides outside the maintained positive model, or changes
of the primal forest. Such moves may be useful, but are different operations.

\begin{lemma}[Monotonicity and tree extension]\label{lem:monotone}
If row deletion set $F$ permits complete peeling, so does every superset of
$F$. Consequently, a successful graphic-independent deletion set can be
extended to a successful spanning tree of $D$.
\end{lemma}
\begin{proof}
Follow any successful peeling order, omitting rows already deleted by the
larger set. A formerly private unit column remains private whenever its
pivot row is still present. A column whose pivot row has already disappeared
can simply be omitted. Equivalently, deleting more rows cannot leave an
original stopping set intact. Finally every forest in a connected graph
extends to a spanning tree, including in a multigraph with loops excluded
from the forest.
\end{proof}

\subsection{The geometric interface is an additional hypothesis}
For the inspected source charts, an extreme edge-gap two-handle has two
authenticated dual endpoint occurrences. A source interior vertex contributes
a three-handle; original-boundary and cut-boundary endpoints are represented
by the root. Contracting a rooted spanning tree gives legal two/three
cancellations and preserves the lower zero/one body and bundle port charts
\cite{dual}. A repeated interior endpoint is two intersections, not one.

Our abstract selector only proposes which such rows to delete. Its graph
input does not prove that these endpoints came from an actual source. The
maintained checker reconstructs endpoints independently from vertex-link
triangles and normal-coordinate edge weights. A proposed native record must
pass that checker, followed by the ordinary core, port, boundary-pattern,
and whole-cut checks. The article's main theorem is complete for the stated
fixed matrix/graph problem, not for all possible downstream geometric
objectives.

\section{Activation closure and the small seed lemma}\label{sec:activation}
\subsection{Low-column implications}
Call a column \emph{usable} if it has at least one entry equal to one.
Columns without a unit entry can never be peeling pivots. Build a directed
graph $G$ on the row set $E$ using usable columns of support at most two.
A unit singleton is an \emph{anchor}; let $A$ be the set of anchors. For a
column with support $\{u,v\}$, include the arc $v\to u$ when $n_{uh}=1$,
and include $u\to v$ when $n_{vh}=1$. In particular $(1,2)$ gives a directed
implication from the heavy row to the unit row, not an undirected edge.

For $S\subseteq E$, let $\Reach(S)$ include $S$ and all vertices reachable
from it in $G$. Low-column peeling from seed set $F$ deletes exactly
$\Reach(A\cup F)$. The forward implication follows by reading a peeling
step as an arc; the reverse follows by following directed paths. Once a
column is used, all rows in its support are already deleted, so retaining
its implication for closure purposes introduces no spurious later move.

Let $\mathcal H$ be the usable columns with support at least three, and define
\begin{equation}\label{eq:B}
 B=\bigcup_{h\in\mathcal H}\supp(h),\qquad b=|B|.
\end{equation}
All support rows, including entries at least two, belong to $B$. Omitting
heavy prerequisites is not a valid reduction. By contrast, a high-support
column with no unit entry contributes nothing to $B$.

A source strongly connected component of $G$ is \emph{mandatory} if it
contains neither an anchor nor a row of $B$. Write these components as
$U_1,\ldots,U_q$. They are pairwise disjoint. Define
\begin{equation}\label{eq:R0}
 R_0=\Reach\left(A\cup\bigcup_{i=1}^{q}U_i\right).
\end{equation}
The use of entire $U_i$ in this symbolic baseline does not grant their rows
free deletions. It is justified only after a genuine seed is chosen in each
$U_i$. Any one row of a strongly connected component reaches all of it and
has the same low-column descendants as every other row in the component.

\begin{lemma}[Mandatory representatives]\label{lem:mandatory}
Every successful initial deletion set meets each $U_i$. If it contains one
row of every $U_i$, then all of $R_0$ can be deleted using anchors and
low-column implications.
\end{lemma}
\begin{proof}
Consider the first row of an unseeded $U_i$ that would be peeled. It cannot
be an anchor. A low-column predecessor outside $U_i$ would be an incoming
arc to a source component, and a high-column pivot would put its row in $B$.
Both are impossible. Conversely, one seed propagates through its strongly
connected component and all its descendants.
\end{proof}

\subsection{A closure operator only on the exceptional boundary}
For each row $e$ define its influence mask
\[
 H(e)=B\cap\Reach(\{e\}).
\]
For $S\subseteq B$, define $\cl(S)$ as the smallest subset of $B$ containing
$S\cup(B\cap R_0)$ and closed under the following two sorts of implications:
\begin{align}
 x\in\cl(S)&\ \Longrightarrow\ H(x)\subseteq\cl(S)
                  &&(x\in B),\label{eq:low-rule}\\
 \supp(h)\setminus\{x\}\subseteq\cl(S)
       &\ \Longrightarrow\ x\in\cl(S)
                  &&(h\in\mathcal H,\ n_{xh}=1).\label{eq:high-rule}
\end{align}
Every prerequisite of \eqref{eq:high-rule} is in $B$, so this is a finite
closure system on $b$ bits. It is extensive, monotone, and idempotent. The
baseline $\cl(\varnothing)$ may contain high-column consequences as well
as $B\cap R_0$. A direct saturation takes polynomial time per mask, and
at most $2^b$ masks ever need memoization.

\begin{theorem}[Exact activation reduction]\label{thm:closure}
Let $I$ consist of one representative from each mandatory component, and let
$X\subseteq E$ be any additional seed set. The exceptional rows deleted by
full peeling from $I\cup X$ are exactly
\begin{equation}\label{eq:activation}
 \cl\left(\bigcup_{e\in X}H(e)\right).
\end{equation}
In particular full peeling succeeds if and only if this set is $B$.
\end{theorem}
\begin{proof}
Each generating implication is a valid low-column path or a valid private
high-column pivot. Lemma~\ref{lem:mandatory} supplies the baseline. Thus the
right-hand side of \eqref{eq:activation} is attained by genuine peeling.

For the converse, inspect a full peeling sequence. A seed in $X$ supplies
$H(e)$, and the mandatory seeds and anchors supply $R_0$. Until a high pivot
occurs, all newly deleted rows lie in the corresponding low reachability
closure. The target and every prerequisite of a high pivot are in $B$;
therefore \eqref{eq:high-rule} already accounts for it. Low propagation from
that target is accounted for by \eqref{eq:low-rule}. Induction along the
sequence proves equality of the exceptional projections.

If all rows are deleted then certainly all of $B$ is deleted. Conversely,
every source component of $G$ either has an anchor, is mandatory, or meets
$B$. Once $B$ is deleted, each source component has a deleted vertex. Every
component of the finite condensation graph lies downstream of a source.
Low-column propagation therefore deletes all rows, including rows outside
$B$.
\end{proof}

\begin{theorem}[At most $b$ useful additional seeds]\label{thm:few-seeds}
A successful spanning tree exists if and only if there are a set $X$ with
$|X|\leq b$ and representatives $u_i\in U_i$ such that
\begin{equation}\label{eq:small-witness}
 X\cup\{u_1,\ldots,u_q\}\text{ is a forest},\qquad
 \cl\left(\bigcup_{e\in X}H(e)\right)=B.
\end{equation}
The representatives and $X$ can be chosen disjoint. Furthermore $X$ has an
ordering in which every added seed strictly increases the closed state.
\end{theorem}
\begin{proof}
Take a successful tree $T$. Choose one representative $u_i\in T\cap U_i$
for every mandatory component, which is possible by
Lemma~\ref{lem:mandatory}. Start at $S=\cl(\varnothing)$ and repeatedly
choose an edge $e$ of $T$ for which $\cl(S\cup H(e))\ne S$.

If this process stopped at a proper subset of $B$, every $H(e)$ for $e\in T$
would already be contained in $S$. Idempotence would then give
$\cl(\bigcup_{e\in T}H(e))=S$, contradicting successful peeling and
Theorem~\ref{thm:closure}. Every chosen edge adds at least one previously
absent bit, so at most $b$ edges are chosen. No row in a mandatory component
can be chosen: its influence is already contained in $B\cap R_0$. Hence
$X$ and the representatives are disjoint, and their union is a subset of
$T$, thus a forest.

Conversely, \eqref{eq:small-witness} and Theorem~\ref{thm:closure} give a
successful forest deletion set. Lemma~\ref{lem:monotone} extends it to a
successful tree. The greedy construction gives the claimed strict-progress
ordering. More generally, a seed whose influence is already in the current
closed state can be omitted permanently: closure idempotence makes it
irrelevant to all later activation, and removing an edge cannot hurt graphic
independence.
\end{proof}

This theorem is the point at which full exceptional-row orders disappear.
We need neither an order of all $b$ rows nor a partition of demands sharing a
seed. A single seed's entire influence mask is inserted at once.

\section{Mandatory choices as a generic contraction}\label{sec:generic}
\subsection{The graph representation}
Orient the nonloop edges of $D$ arbitrarily. Let
$V_D=(v_e)_{e\in E}$ be its reduced incidence matrix over a prime field
$\F_p$, with the row of the boundary root removed. It has $r$ rows.
A loop has the zero column. Over every field, the columns indexed by a set
of edges are linearly independent if and only if that edge set is a forest.
For a forest, repeatedly removing a nonroot leaf produces a square minor
with determinant $1$ or $-1$. A cycle, including a loop or a two-edge
parallel cycle, gives a column dependence. This also explains why
characteristic two is harmless for the \emph{graph representation}.

Introduce an independent indeterminate $z_{i,e}$ for each $e\in U_i$, and set
\begin{equation}\label{eq:generic-y}
 y_i=\sum_{e\in U_i}z_{i,e}v_e,\qquad
 Y=[y_1\ \cdots\ y_q].
\end{equation}
These columns initially live over the rational function field
$\F_p(\mathbf z)$.

\begin{lemma}[Generic representative contraction]\label{lem:generic}
For any set $X$ of $k$ edge labels, the following are equivalent:
\begin{enumerate}[label=(\roman*)]
\item $[Y,V_{D,X}]$ has rank $q+k$ over $\F_p(\mathbf z)$;
\item one can choose representatives $u_i\in U_i$, distinct from one another
and from $X$, for which $X\cup\{u_i:1\leq i\leq q\}$ is a forest.
\end{enumerate}
If a specialization of the $z_{i,e}$ has rank $q+k$, condition (ii) holds
for the original graph, without any probabilistic qualification.
\end{lemma}
\begin{proof}
Expand any $(q+k)$-row minor in the first $q$ columns. Its terms are indexed
by tuples $(u_1,\ldots,u_q)$, with monomial
$\prod_i z_{i,u_i}$ and coefficient equal to the corresponding graph-incidence
minor with the columns $X$ appended. Different tuples give different
monomials because the variables retain both the color and the edge label.
Thus a minor is a nonzero polynomial precisely when at least one of these
coefficients is nonzero.

A nonzero coefficient implies a forest of $q+k$ distinct edges. Conversely,
a forest of representatives together with $X$ has a graph minor equal to
$\pm1$. Its monomial coefficient cannot cancel with any other tuple, and
remains nonzero over every prime field. Finally, a polynomial that evaluates
to a nonzero value is not the zero polynomial. Specialization may destroy
independence but cannot create an independent set that lacks genuine
representatives.
\end{proof}

If no independent representatives exist even for $X=\varnothing$, the local
problem is negative. A polynomial graphic/partition matroid-intersection
prepass detects this and can provide a Rado obstruction: a subcollection
$J\subseteq\{1,\ldots,q\}$ with
\begin{equation}\label{eq:rado}
 |J|>\rk_D\left(\bigcup_{i\in J}U_i\right).
\end{equation}
The delivered checker verifies this inequality directly from graph
components. This classical prepass is an optimization and an exact-negative
route; it is not the new single-exponential argument.

\subsection{Specialize, quotient, and reduce the rank}
Choose every $z_{i,e}$ uniformly and independently in $\F_p$. If the sampled
$Y$ has rank less than $q$, abandon this round. Contracting fewer than $q$
independent columns would forget mandatory constraints and is not allowed.
Otherwise compute a linear quotient map
\[
 Q:\F_p^r\longrightarrow\F_p^{r-q},\qquad\ker Q=\langle Y\rangle.
\]
Gaussian elimination constructs $Q$: reduce a vector by a normalized echelon
basis of $Y$ and retain its free coordinates. Put
$d=\min\{b,r-q\}$ and independently choose a uniform random matrix
$P\in\F_p^{d\times(r-q)}$. The small representation used by the search is
\begin{equation}\label{eq:small-rep}
 a_e=P Q v_e\in\F_p^d.
\end{equation}
A set independent in this representation was independent in the quotient
before projection, and therefore has genuine mandatory representatives by
Lemma~\ref{lem:generic}. Rank reduction only loses potential witnesses.

\begin{remark}[Do not apply the field to geometric counts]
The field acts on $V_D$, not on $N$. For example,
$\left(\begin{smallmatrix}2&1\\1&2\end{smallmatrix}\right)$ is a stopping
matrix, but its reduction modulo two is a permutation matrix and peels
completely. Replacing unsigned incidences by residues would change the
problem. The implementation keeps all heavy prerequisites even in the
fully deterministic characteristic-two search.
\end{remark}

\subsection{A witness can be unmixed deterministically}
\begin{proposition}[Algebraic unmixing]\label{prop:unmix}
Suppose sampled $Y$ and a set $X$ have $\rk[Y,V_{D,X}]=q+|X|$.
Actual independent representatives can be recovered in polynomial time
using only field-rank tests and the lists $U_i$.
\end{proposition}
\begin{proof}
Maintain the independent list consisting of the $q$ current columns and
$V_{D,X}$. At step $i$, its $i$th column is still the original sampled
linear combination $y_i\in\langle v_e:e\in U_i\rangle$. Its exterior
product with the other columns is nonzero. Multilinearity implies that
replacing $y_i$ by at least one $v_e$, $e\in U_i$, leaves a nonzero exterior
product. Scan the candidates with rank tests, choose such an edge, and keep
that replacement. After $q$ steps every mandatory column is an actual edge.
Independence guarantees that no edge was repeated or already belonged to
$X$. The number and dimensions of the rank tests are polynomial in the
explicit input size.
\end{proof}
The current producer instead invokes its already tested graphic/partition
oracle with $X$ contracted, to recover representatives. Proposition~\ref{prop:unmix}
is a proved alternative, not an additional executed recovery path in the
package. It also shows that mandatory matroid intersection is not logically
necessary for the randomized positive-search theorem: the exact rank-
obstruction prepass may be omitted and all positive witnesses can be recovered
by linear algebra.

\section{Exterior representative families over closed states}\label{sec:dp}
\subsection{The representative property}
For an ordered independent $k$-tuple $X=(e_1,\ldots,e_k)$ in the small
representation, write
\[
 w_X=a_{e_1}\wedge\cdots\wedge a_{e_k}
     \in\bigwedge^k\F_p^d.
\]
This space has dimension $D_k=\binom dk$. The coordinates of $w_X$ are the
$k\times k$ minors of the matrix with columns $a_{e_i}$. In particular,
$w_X\ne0$ exactly when $X$ is independent.

\begin{lemma}[A spanning subfamily preserves extensions]\label{lem:rep}
Let $\mathcal F$ be a family of independent $k$-sets having the same closed
activation state $S$. Choose actual members $\widehat{\mathcal F}\subseteq
\mathcal F$ whose exterior vectors form a basis of the span of all
$\{w_X:X\in\mathcal F\}$. Then
$|\widehat{\mathcal F}|\leq\binom dk$. For every future edge set $Z$ with
$|Z|\leq d-k$, if some $X\in\mathcal F$ has $X\cup Z$ independent,
then some $\widehat X\in\widehat{\mathcal F}$ does too. Starting from
$X$ or $\widehat X$, the future activation caused by the same $Z$ is identical.
\end{lemma}
\begin{proof}
The map $w\mapsto w\wedge w_Z$ is linear. If it is nonzero on any vector
in the span, it is nonzero on at least one vector of a spanning basis.
A nonzero product guarantees both independence and disjointness from $Z$.
The activation state records everything a future seed can affect:
\[
 \cl\left(S\cup\bigcup_{e\in Z}H(e)\right).
\]
It does not depend on the particular set used to reach $S$.
\end{proof}
This is the standard linear-matroid representative-family mechanism
\cite{fomin}. Its application here is valid because
Theorem~\ref{thm:closure} has separated all future activation information
from the identity of the already selected seed edges.

\subsection{The dynamic program}
At state $(S,k)$ retain a basis of exterior vectors of actual $k$-seed
witnesses whose closed state is $S$. Start with the empty witness and the
scalar exterior vector $1$ at $(\cl(\varnothing),0)$. For every retained
witness $X$ and every edge $e$, form
\begin{equation}\label{eq:transition}
 S'=\cl(S\cup H(e)),\qquad w'=w_X\wedge a_e.
\end{equation}
Ignore the transition if $S'=S$ or $w'=0$. Otherwise insert its actual
witness $X\cup\{e\}$ into the basis at $(S',k+1)$ if its exterior vector
is not already in that span. If $S'=B$, return the witness immediately.
Process states by increasing $k$; only $0\leq k\leq d$ can occur.

\begin{theorem}[Completeness in a supplied representation]\label{thm:dp}
For any fixed vectors $a_e\in\F_p^d$, the dynamic program returns a witness
if and only if some independent edge set $X$ satisfies
$\cl(\bigcup_{e\in X}H(e))=B$.
\end{theorem}
\begin{proof}
Every retained witness is actual, independent, and has its recorded closed
state, by construction. This proves soundness.

For completeness, choose a successful independent witness and remove
redundant seeds as in Theorem~\ref{thm:few-seeds}. Its seeds can be added in
strict-progress order. When an intermediate witness is discarded in favor
of a basis, Lemma~\ref{lem:rep} preserves an extension by the remaining
ordered seed set. The replacement has the same state, and the nonzero
exterior product with all remaining seeds guarantees that every prefix of
that extension remains independent. The next activation state is therefore
unchanged. Induct on the number of remaining additions until the full
boundary is reached. This argument also covers replacements of a witness
that already arose as a previous representative replacement.
\end{proof}

\subsection{A representation invariant that matters in code}
Gaussian elimination of a candidate exterior vector can produce a linear
combination that is not the exterior vector of \emph{any} actual seed set.
Consequently an eliminated basis row must not be treated as a new witness.
The implementation stores two objects separately: normalized echelon rows
for span tests, and each retained original decomposable vector together with
its actual ordered edge tuple. Transitions extend only the latter. A dedicated
regression test distinguishes the raw vector from its eliminated row.

The literal-set ablation also retains actual witnesses. It removes repeated
\emph{sets} reached through different orders, but does not apply exterior
span compression. The measured saving is therefore not merely suppression
of different permutations of an identical edge set.

\section{The single-exponential search theorem}\label{sec:main}
\subsection{Why only one witness has to survive randomness}
We use the elementary polynomial zero bound: a nonzero polynomial of total
degree at most $t$ over $\F_p$, evaluated at independent uniform field
values, vanishes with probability at most $t/p$. For completeness, write it
as a polynomial in its last variable, with leading coefficient of degree
at most $t-s$ and last-variable degree $s$. The leading coefficient vanishes
with probability at most $(t-s)/p$ by induction; otherwise there are at most
$s$ roots in the last coordinate. This proves the bound. It is the usual
Schwartz--Zippel argument \cite{schwartz}.

\begin{lemma}[Fixed-witness survival]\label{lem:survival}
Fix a feasible $k$-seed witness $X_*$ in
Theorem~\ref{thm:few-seeds}. One random specialization and projection retain
its independence with probability at least
\begin{equation}\label{eq:miss}
 1-\frac{q+k}{p}.
\end{equation}
The witness is fixed before the random variables of this round are sampled.
\end{lemma}
\begin{proof}
By Lemma~\ref{lem:generic}, some $(q+k)$-minor of
$[Y,V_{D,X_*}]$ is a nonzero polynomial. Its degree is at most $q$, because
only the mandatory columns are random. Thus this augmented matrix loses
rank with probability at most $q/p$. On the complementary event $Y$ has
rank $q$, so the round is not abandoned, and the quotient images of $X_*$
are independent.

Condition on any such specialization of $Y$. The independent quotient
images have a $k$-dimensional span. Since $k\leq d$, there exists a linear
map to $\F_p^d$ that preserves them. Accordingly some $k$-minor of the
random projected columns is a nonzero polynomial of degree at most $k$ in
the independent entries of $P$. Its failure probability is at most $k/p$.
The conditional bound holds for every successful first-stage specialization.
Adding the two failure probabilities proves \eqref{eq:miss}.
\end{proof}

No union bound over every feasible seed set is needed. On a positive input,
one fixed successful witness suffices. Once the small representation retains
that witness, Theorem~\ref{thm:dp} completes its search deterministically.
Requiring a simultaneously faithful representation of all $m^{O(b)}$
possible seed sets would be an unnecessary extra demand.

\begin{theorem}[One-sided single-exponential selector]\label{thm:main}
For the fixed local problem of Section~\ref{sec:model}, there is an algorithm
with these properties:
\begin{enumerate}[label=(\roman*)]
\item Every reported success has a valid spanning tree and complete peeling
certificate, checked deterministically from the original graph and matrix.
\item If a successful tree exists, $\lambda$ uncapped independent rounds
find one with probability at least $1-2^{-\lambda}$.
\item Its bit running time is $2^{O(b)}\poly(\Lambda,\lambda)$.
One round admits \eqref{eq:headline} as a conservative field-operation bound,
with $d=\min\{b,r-q\}$ when mandatory representatives are feasible.
\item Its field-element storage can be bounded by
$2^b4^d\poly(\Lambda)\leq8^b\poly(\Lambda)$.
\end{enumerate}
On an unsuccessful search the general algorithm reports \Unk, not a
negative topological answer. An optional exact rank-obstruction route may
certify local infeasibility sooner.
\end{theorem}
\begin{proof}
First handle the empty input and any automatically complete closure. The
mandatory prepass either gives the verified obstruction
\eqref{eq:rado}, or ensures $q\leq r$. Alternatively it can be omitted and
rank-deficient mandatory samples can simply be abandoned.

Choose a prime $p>2(q+b)$ with $p=O(q+b+1)$; elementary prime bounds and
trial division suffice in polynomial time for these explicitly bounded
parameters. The implementation takes the first prime at least
$\max\{101,2(q+b)+1\}$. Sample $Y$, check its rank, quotient, and project
as in Section~\ref{sec:generic}. Search by Theorem~\ref{thm:dp}.
If $X$ is found, sampled independence implies genuine extendibility by
Lemma~\ref{lem:generic}. Recover mandatory representatives by the
polynomial forest oracle or Proposition~\ref{prop:unmix}; extend their
union with $X$ to a tree. Then replay the actual matrix peeling certificate.
All these steps are deterministic after the sample, and no field outcome
can cause a false positive.

On a positive input fix the witness furnished by
Theorem~\ref{thm:few-seeds}, independently of the sampling. By
Lemma~\ref{lem:survival} a round misses it with probability less than
$1/2$. Conditional on retaining it, the dynamic program succeeds. Independent
repetition gives the stated miss bound.

For the complexity, each mask and cardinality state retains at most
$D_k=\binom dk$ witnesses. There are at most $2^b$ masks. Therefore the
number of candidate edge extensions is at most
\[
 m\,2^b\sum_{k=0}^{d}D_k=m\,2^{b+d}.
\]
An exterior extension costs $O(dD_k)$ arithmetic operations in a direct
coordinate implementation. A straightforward insertion into a basis in
dimension $D_{k+1}$ costs $O(D_{k+1}^2)$. Since
\[
 \sum_{k=0}^{d-1}D_kD_{k+1}^2
 \leq \left(\sum_kD_k\right)\left(\max_kD_k\right)^2
 \leq8^d,
\]
the claimed time bound follows. Echelon bases and raw witness vectors occupy
at most a polynomial factor times
\[
 2^b\sum_{k=0}^{d}\binom dk^2
   =2^b\binom{2d}{d}\leq2^b4^d
\]
field elements. Influence preparation, closure saturation, graph operations,
quotients, and certificate output are polynomial factors in $\Lambda$.
Field elements have $O(\log(q+b+2))$ bits. Standard exact finite-field
arithmetic gives the stated bit bound. The research implementation uses
sparse dictionaries; even conservative dictionary overhead is absorbed in
$2^{O(b)}\poly(\Lambda,\lambda)$, rather than being asserted to have a
particular machine-level base of sixteen.
\end{proof}

The algorithm's probability statement assumes fresh independent uniform
samples and enough resources to complete the claimed rounds. A fixed
pseudorandom seed is useful for regression reproduction, not a derandomization.
A boundary or transition cap may cause \Unk\ even on a positive input, so a
capped call is not automatically entitled to the uncapped miss guarantee.

\section{Deterministic low-rank decisions and frontier certificates}\label{sec:frontier}
\subsection{A class with no randomness at all}
\begin{theorem}[Faithful low-rank selection]\label{thm:faithful}
Suppose $q=0$ and $r\leq b$. Successful dual-tree selection is decidable
deterministically in $2^{O(b)}\poly(\Lambda)$ time. A positive answer has
the same tree-and-peeling certificate as before. A negative answer has a
frontier-span certificate of at most $4^b\poly(\Lambda)$ bits, verifiable
independently in single-exponential time.
\end{theorem}
\begin{proof}
Use the original reduced binary incidence columns $v_e\in\F_2^r$.
They represent graphic independence exactly; no mandatory contraction or
rank projection is necessary. Run the dynamic program with $d=r$. Its
soundness and completeness follow from Theorem~\ref{thm:dp} and the small
seed theorem. A negative answer can retain the successful basis insertions
at every state. There are at most $2^b\sum_k\binom rk\leq2^{b+r}$
actual witnesses, each with at most $r$ edge labels. The certificate and
verification claims follow from the explicit construction below.
\end{proof}

\subsection{What the negative packet contains}
A frontier packet contains an input digest and, for every retained state
$(S,k)$, a list of actual $k$-edge seed witnesses whose exterior vectors are
a basis for that state. It contains no alleged field-operation transcript
and no claim that a reduced linear combination is an actual witness.
The complete boundary state $B$ must not occur.

An independent verifier performs the following checks.
\begin{enumerate}
\item Recompute $B$ and the mandatory components from the supplied matrix;
verify $q=0$ and $r\leq b$. Check input binding, unique state keys, label
ranges, and the absence of a full-boundary state.
\item For each alleged witness $X$, recompute its exterior coordinates as
binary determinants of graph-incidence minors. Check that the vector is
nonzero and that literal matrix peeling from $X$ has precisely the claimed
exceptional projection. Verify that each state's listed vectors are
linearly independent.
\item Check that the empty exterior vector $1$ is in the initial state's
span. For each state witness $X$ and each edge $e$ giving strict activation
progress, recompute $w_{X\cup\{e\}}$. If nonzero, require it to belong
to the listed span at the appropriate next state. At $k=r$ every further
edge extension is dependent, so there is nothing to check.
\end{enumerate}
The supplied implementation computes activation in this checker by literal
peeling, not by importing the producer's Horn-closure code. It computes
minors directly, not by importing the producer's wedge-extension routine.

\begin{lemma}[Frontier coverage proves local nonexistence]\label{lem:frontier}
If these checks pass, no successful tree exists for the supplied local input.
\end{lemma}
\begin{proof}
The initial state spans the empty witness. Induct on a strict-progress seed
sequence. Suppose the exterior vector of its current prefix is in the span
of the listed state basis. Wedge multiplication by the next edge is linear,
and its next activation state depends only on the current state and that
edge. The transition checks therefore place the next prefix vector in the
listed next-state span. If a successful independent seed set existed, choose
a strict-progress ordering as in Theorem~\ref{thm:few-seeds}. Its final
vector would be nonzero in a full-boundary state, contrary to the packet's
absence of such a state. A successful tree supplies such a seed set.
\end{proof}

This is an exact local negative certificate, unlike a failed random
projection. It must not be extended silently to $q>0$, where a specialized
mandatory contraction can lose witnesses, or to $r>b$, where a sampled
projection may do so. The verifier explicitly rejects those unsupported
regimes. Nor is a valid frontier packet a certificate of knottedness.

\section{Necessary parity checks for geometric source data}\label{sec:parity}
The abstract selector deliberately allows matrices and graphs that may not
come from a three-manifold. There is nevertheless a useful inexpensive
source-consistency condition, separate from the selection theorem.

\begin{proposition}[Dual-cycle parity]\label{prop:parity}
Suppose $N$ gives \emph{exact} unsigned one/two-handle intersection counts,
and $D$ gives the actual two/three-handle endpoint occurrences, in a handle
chain complex with the stated lower forest contraction. Then
\begin{equation}\label{eq:parity}
 V_D (N\bmod2)=0\quad\text{over }\F_2.
\end{equation}
Equivalently, for each column of $N$, the dual edges on which the count is
odd have even degree at every dual vertex.
\end{proposition}
\begin{proof}
Over $\F_2$ signs disappear. In handle bases, $\partial_2$ is the transpose
of the unsigned incidence matrix modulo two, and $\partial_3$ is the
transpose of the reduced dual incidence matrix. The chain identity
$\partial_2\partial_3=0$ gives \eqref{eq:parity} after transposing.
The equation at the omitted boundary root follows from the sum of the
other incidence rows. A loop contributes twice and therefore contributes
zero modulo two.
\end{proof}
This is a specialization of the chain identity, not a new general theorem
about chain complexes. It is relevant because it detects a class of abstract
fixtures that must not be advertised as genuine geometric sources.

The test cannot be applied directly to entries that were merely capped at
two. An entry recorded as ``at least two'' may have either actual parity.
There is an exact polynomial criterion for whether the \emph{saturated
categories} admit any compatible parity lift.

\begin{theorem}[Parity lifting for saturated columns]\label{thm:parity-lift}
For one saturated column, let $U$ be its unit edges and $H$ its heavy edges.
Let $O$ be the vertices of odd degree in the dual subgraph with edge set
$U$. A choice of parities for the edges of $H$ makes $U$ together with the
odd heavy edges Eulerian if and only if
\begin{equation}\label{eq:even-demand}
 |O\cap K|\text{ is even for every connected component }K\text{ of }D[H].
\end{equation}
When feasible, a spanning-forest leaf sweep constructs the lift. All heavy
entries can be realized as either $2$ or $3$ for this parity purpose.
\end{theorem}
\begin{proof}
A heavy-edge subgraph has an even number of odd-degree vertices in each
heavy connected component, proving necessity. For sufficiency choose a
spanning forest of $D[H]$. Root each component. Process its nonroot vertices
from leaves toward the root; include a parent edge precisely when its child
currently has odd demand, then toggle that demand at the parent. Only the
root can retain demand, and \eqref{eq:even-demand} makes its final demand
even. The chosen heavy edges have boundary $O$, so their symmetric sum with
$U$ is Eulerian. Assign count three to chosen heavy edges and count two to
the others. Loops need not be chosen.
\end{proof}
The implementation supplies both diagnostics in
\path{integration/parity_audit.py}. Neither authenticates source charts or
constructs a manifold. A parity-compatible lift is not permission to change
an already authenticated exact count. The maintained five-row matrix passes
the exact parity test. The nonlinear choice family below also passes it;
its three-dimensional realization remains unproved.

\section{The maintained loop example and the integration seam}\label{sec:native}
For the source with one interior dual vertex, the maintained matrix on
columns with original labels $(3,5,6,7,8)$ is \cite{loop}
\[
 N=\begin{pmatrix}
 0&0&1&0&0\\
 1&1&0&0&0\\
 0&1&0&1&1\\
 1&0&1&1&1\\
 0&1&1&1&0
 \end{pmatrix}.
\]
Rows $0,1,2,3$ are parallel root-to-interior edges, while row $4$ is an
interior loop. Deleting row $0$ leaves totals $(2,3,2,3,2)$, hence a stopping
set. Here $b=5$, $q=0$, and $r=1$, so Theorem~\ref{thm:faithful} applies.
The new search selects row $1$ after two candidate transitions. Its local
peeling schedule, in normalized column labels, is
\[
 (3,0),\quad(2,4),\quad(4,1),\quad(0,2).
\]
In original column labels this is
\[
 (3,3),\quad(2,8),\quad(4,5),\quad(0,6).
\]
The source's previously recorded alternative schedule used column $7$ at
the third step rather than column $5$. Both are legal matrix schedules.
The new local schedule consequently leaves a different column representative;
one must not copy stale downstream port words or a prior basis-dependent
certificate on the assumption that the chosen tree fixes the entire proof.

The delivered adapter translates only the verified tree into a native
two/three reduction proposal. For a normalized interior identifier $37$,
its cancellation list is the single record
\begin{quote}\ttfamily
\{two\_handle: 1, three\_handle: 37\}.
\end{quote}
It preserves the complete endpoint list, including the repeated endpoint
of row $4$, and records surviving source row IDs $(0,2,3,4)$. Native replay
is responsible for renumbering those survivors. The adapter checks that
its input graph really matches the supplied endpoint records, so the local
certificate cannot silently be applied to another graph.

The maintained target is
\path{fastunknot.core_three_handle_reduction_verify.replay_core_three_handles}.
The required schema is \path{core-three-handle-reduction-v1}, with fields
\path{dual_endpoints}, \path{cancellations}, and
\path{surviving_two_handles}. Its current source checker was read at the
pinned commit. The adapter's record normalization, source-ID mapping,
mismatch refusals, and proposal were executed in seven tests. Full native
geometry, Regina cuts, bundle-port replay, and end-to-end unknot recognition
were \emph{not} executed in this environment. These remain mandatory
integration tests, not results inferred from the local example.

\section{Implementation, resource semantics, and certificates}\label{sec:implementation}
\subsection{Modules and independent paths}
The new package uses only the Python standard library. It was executed with
CPython 3.13.5 on the Linux environment recorded in the result files.
\path{model.py} validates the explicit graph and unsigned matrix;
\path{closure.py} builds reachability and the exceptional closure;
\path{matroid.py} implements a simple disjoint-color graphic/partition
oracle; \path{algebra.py} performs finite-field and exterior operations;
\path{solve.py} orchestrates the routes; and \path{verify.py} imports none
of the producer, graph optimizer, closure, or exterior modules.

The polynomial preprocessing is intentionally straightforward. Low
reachability uses bit masks and transitive closure rather than an optimized
SCC/condensation implementation. The matrix producer uses a dense deterministic
peeling scan. Neither choice is claimed to improve the maintained sparse
peeler; the research target is the exponential \emph{selection} dependence.
The small-parameter method belongs behind a cheap incumbent attempt, as in
the supplied default.

The positive certificate consists of an input digest, a spanning-tree edge
list, its boundary-outward dual order, and a complete peeling schedule. A
simple independent checker verifies the tree, dual moves, private columns,
and full coverage. No random seed, field modulus, Gaussian transcript, or
representative-family calculation is trusted to establish the answer.
A SHA-256 digest prevents accidental input substitution; it is not a proof
that an input came from a triangulation.

\subsection{The three statuses}
The public interface has three meanings. \Yes\ provides a verified local
positive certificate. \No\ provides either a verified mandatory-rank
obstruction or a verified faithful low-rank frontier. \Unk\ means that no
such certificate was obtained within the supported route and resources.
In particular a failed sample, an exhausted round allowance, an exceptional
boundary cap, or a transition cap does not become a local negative answer.
Even \No\ is not a topological negative verdict.

The default caps are $b\leq20$ and one million considered transitions across
sampling rounds, with eight requested rounds. They are practical research
safeguards, not the hypotheses of an unconditional guarantee for every
positive input. Passing \texttt{None} as the corresponding API arguments
runs the uncapped algorithm in the theorem. The callback propagates
interruption exceptions; it does not translate them to a verdict.
Source preparation and polynomial preprocessing are not bounded merely by
a transition counter.

The validator rejects an impossible gigantic vertex count before allocating
vertex arrays, using $|V|\leq m+1$ for a connected explicitly listed graph.
Empty-row input does not expand a potentially large declared count of zero
columns. A $2^{16384}$ matrix coefficient is transported in hexadecimal and
bound exactly by the input digest without decimal-expansion limits. The
capability is tested, not presented as a benchmark of native normal-sheet
multiplicity processing.

\subsection{Reproduction commands}
From the package root:
\begin{verbatim}
python -m unittest discover -s tests -v
python -m experiments.validate
python -m experiments.benchmark
python -m activation_forest examples/repository_loop.json --no-incumbent
python -m activation_forest examples/generic_positive.json --seed 1210
python -m activation_forest examples/repository_loop.json \
  --verify results/repository_loop_answer.json
cd article && sh build.sh
\end{verbatim}
Fixed seeds in the archived experiments make their finite executions
reproducible. Production use without \texttt{--seed} requests fresh system
randomness. The old selector is vendored unchanged with its license and
provenance solely for comparative reproduction; it is not counted as new
implementation code.

\section{Executed validation and measured tradeoffs}\label{sec:experiments}
\subsection{A ledger that distinguishes unknowns from negatives}
The tests and audit scripts were executed, rather than only supplied.
\texttt{results/validation.json} contains environment details, counters,
seeds, and the driver hash. The following table summarizes distinct audit
loops; unit-test assertions are not silently added to the input-case count.

% BEGIN VALIDATION_TABLE
\begin{center}\small
\begin{tabular}{lrrrr}\toprule
Audit & Trials & Positive & Exact negative & Unknown\\\midrule
Exhaustive $4	imes2$ & 6,561 & 2,586 & 2,555 & 1,420\\
Random matrix/graph & 1,500 & 531 & 836 & 133\\
Planted generic positives & 300 & 300 & 0 & 0\\
\bottomrule\end{tabular}
\end{center}
\begin{center}\small\begin{tabular}{lr}\toprule
Additional algebraic audit & Queries or samples\\\midrule
Forced matroid extensions & 16,038\\
Horn closure versus literal peeling & 2,369\\
Sampled-independence soundness & 1,951\\
Projected representative DP versus literal sets & 255\\
Complete small-field sample spaces & 503\\
\bottomrule\end{tabular}\end{center}
% END VALIDATION_TABLE

The exhaustive matrix audit fixes a three-vertex four-edge dual multigraph
with a parallel edge and visits every one of the $3^8=6561$ four-by-two
matrices with entries in $\{0,1,2\}$. The random audit varies connected
multigraphs, including loops, and dense matrices. Both compare with literal
enumeration of every spanning tree. All 3,117 positive inputs in these two
sets receive positive certificates. Of their 4,944 negative inputs, 3,391
receive exact local certificates and 1,553 remain \Unk. Those inconclusive
outputs are not counted as negative decisions or as proofs of full local
completeness.

A further 300 planted positive trials use cycle dual graphs, $3\leq b\leq7$,
and $1\leq q\leq6$. All succeed, requiring 318 total sampled rounds.
These abstract families exercise the mandatory contraction and random rank
reduction, rather than only the deterministic $q=0$ route. They are not
asserted to be geometric source matrices.

The 71 unit tests include positive and negative certificate corruption,
producer-disabled replay, bad random samples, resource caps, interruption,
loops and repeated endpoints, large binary counts, the record-level adapter,
and the exact versus saturated parity distinction. The matroid-intersection
audit compares 16,038 forced-extension queries against literal representative
tuples. The projected-search audit compares the new dynamic program with
literal independent-seed enumeration in exactly the same sampled
representation; it does not confuse a sampled miss with an error in that
representation's search.

\subsection{An exact finite-field probability check}
Take the triangle with edge vectors $v_0,v_1,v_2$, one mandatory class
$\{0,1\}$, and the fixed additional witness $\{2\}$. Sampling
$y=z_0v_0+z_1v_1$ and a one-dimensional quotient projection uses three field
samples. The quotient witness survives the first stage with probability
$1-1/p$, and the nonzero scalar projection survives with the same conditional
probability. Thus the exact miss probability is
\begin{equation}\label{eq:small-field}
 1-(1-1/p)^2=\frac{2p-1}{p^2}\leq\frac2p.
\end{equation}
The script enumerates all $503$ samples over $p=2,3,5,7$. The observed miss
counts are respectively $6/8$, $15/27$, $45/125$, and $91/343$, exactly as in
\eqref{eq:small-field}. This is a complete finite sample-space check, not a
statistical estimate. Small fields can have substantial miss probability;
that is why the theorem's constant-success choice is $p>2(q+b)$.

\subsection{A nonlinear family with exact witness-count separation}
For $g\geq1$, let the dual graph have root $0$ and $g$ leaves. On arm $i$
place three parallel edges, in input order $c_i,a_i,b_i$. Add a loop $z$ at
the root. Introduce a column with unit entries on $a_i,b_i$, and a column
with entries $1$ on $c_i$ and $3$ on $a_i$. Finally introduce one global
column having entry $2$ on every $a_i,b_i$ and entry $1$ on $z$. All
unmentioned entries vanish. The use of $3$ is intentional: it has the same
activation category as $2$, but makes this exact-count family satisfy
Proposition~\ref{prop:parity}.

There are $3g+1$ rows and $2g+1$ columns. The exceptional boundary consists
of all $a_i,b_i$ and $z$, so $b=2g+1$, $q=0$, and $r=g$. There is no initial
anchor. Each $a_i,b_i$ is a low strongly connected pair, and deleting either
also deletes $c_i$. The loop $z$ can be deleted only when every good pair
has disappeared, through the genuinely higher-support final column.

\begin{proposition}[Choice-family counts]\label{prop:family}
A spanning tree of this family succeeds exactly when it chooses $a_i$ or
$b_i$ on each arm. There are $2^g$ successful trees out of $3^g$.
Before the first complete state is reached, the representative search retains
$2^g-1$ witness vectors, while deduplicated literal seed-set storage retains
$3^g-2^g$ witnesses. In the inspected predecessor's specified input and
partition order, the first successful exceptional-order branch is its
$\Bell(g)$th matroid-intersection subproblem.
\end{proposition}
\begin{proof}
A tree chooses exactly one of the three edges on each arm and cannot choose
$z$. Choosing $a_i$ or $b_i$ activates the entire good pair and then $c_i$.
Choosing only $c_i$ leaves the good pair unavailable: the global column has
no unit there, and neither low column can delete the first good row. Hence
every arm must have a good choice, after which the global column deletes $z$.
This proves the tree count.

For a proper subset $J$ of activated arms, there are $2^{|J|}$ distinct
independent useful seed sets, one good edge on each arm of $J$. In the
binary star-incidence representation all these sets have the same nonzero
exterior vector. One representative per $J$ therefore suffices, giving
$2^g-1$ representatives. Literal deduplicated storage gives
$\sum_{J\subsetneq[g]}2^{|J|}=3^g-2^g$.

In the predecessor's first exceptional-row ordering the first row of each
good pair imposes a demand with candidate set $\{a_i,b_i\}$; its partner
is automatically reached, and $z$ is reached at the end. The $g$ demand
sets are disjoint, so any partition merging two demands has an empty block
candidate set. Only the discrete partition succeeds, and it is last in
the audited restricted-growth partition order. Thus $\Bell(g)$ calls are
made. This last assertion concerns the archived implementation order, not
an inherent lower bound for every algorithm using exceptional orders.
\end{proof}

The graph rank and closure make this a useful ablation family, not a proposed
hard family of knot diagrams. Satisfying the parity condition does not
establish a three-dimensional realization. A specialized recognizer of its
arm structure could solve it faster than either general-purpose method.

\subsection{Paired local-query timings}
Each timing begins with a validated matrix/graph instance and includes
preprocessing, selection, certificate production, and one independent
replay. The new producer already performs its replay before returning;
a redundant second replay is not charged. The historical producer is
followed by its independent checker. Source geometry construction, file
I/O, and full knot recognition are excluded. Literal positive tree searches
are also replayed. Literal negative enumeration has no separate compressed
negative certificate and is reported as an exact-search control, not as
an identical proof-format workflow.

The table gives median per-query milliseconds and the median of paired
historical/new ratios. Five randomized-order pairs are retained per case.
Submillisecond queries were batched; earlier unbatched data and earlier
anchored-family controls are also retained and explicitly named. No case
was removed merely because it favored the older or literal method.

% BEGIN BENCHMARK_TABLE
\begin{center}\small\setlength{\tabcolsep}{5pt}
\begin{tabular}{lrrrr}\toprule
Local input & Prior (ms) & New (ms) & Prior/new & Literal (ms)\\\midrule
Parallel negative, 3 & 0.1932 & 0.0786 & 2.61 & 0.0036\\
Parallel negative, 5 & 28.8081 & 0.1113 & 259 & 0.0068\\
Parallel negative, 6 & 642.6735 & 0.1358 & 4.78e+03 & 0.0080\\
Maintained loop matrix & 1.1715 & 0.0733 & 17.4 & 0.0366\\
Nonlinear choice, 2 arms & 0.0929 & 0.0868 & 1.07 & 0.0737\\
Nonlinear choice, 5 arms & 1.0083 & 0.6783 & 1.5 & 8.0923\\
Generic positive & 0.0596 & 0.0913 & 0.653 & 0.0202\\
\bottomrule\end{tabular}\end{center}
% END BENCHMARK_TABLE

The reconstructed maintained loop matrix benefits relative to the audited
exceptional-order implementation, but literal testing of its very few trees
is still faster. The generic positive example favors the prior implementation.
The enormous historical/new ratio on the six-parallel-edge negative instance
is particularly unrepresentative of generic difficulty: only six actual trees
exist, and literal enumeration is much faster than either certificate-search
framework. It measures avoidance of needless old order/partition branches,
not evidence that difficult knot instances have been solved thousands of
times faster. Odd-size members of that all-unit parallel family fail even
the exact parity condition; the even-size member passes it, without proving
geometric realizability.

For the nonlinear choice-family ablation, both arms perform the same
activation preparation and discrete positive replay; only the compression
of candidate seed families differs.

% BEGIN ABLATION_TABLE
\begin{center}\small\setlength{\tabcolsep}{5pt}
\begin{tabular}{rrrrrr}\toprule
Arms & Literal witnesses & Representatives & Literal (ms) & New (ms) & Ratio\\\midrule
4 & 65 & 15 & 0.455 & 0.271 & 1.672\\
6 & 665 & 63 & 6.380 & 1.297 & 4.869\\
8 & 6,305 & 255 & 87.074 & 6.453 & 13.494\\
\bottomrule\end{tabular}\end{center}
% END ABLATION_TABLE

At eight arms, $6,305$ deduplicated literal witnesses are reduced to $255$
representatives, as predicted by Proposition~\ref{prop:family}. All samples,
operation counts, input hashes, source hashes, and same-algorithm A/A
controls are in \path{results/benchmark.json}. These are finite Python
measurements, not an empirical determination of an asymptotic exponent.

\section{What this contributes to a quasi-polynomial recognizer}\label{sec:qp}
\begin{corollary}[Conditional local quasi-polynomial regime]\label{cor:qp}
Let a recognition workflow on an $n$-crossing input make at most
$L(n)=2^{O((\log n)^2)}$ calls to this local selector. Suppose every call
that must succeed has a successful forest, its encoded local input size is
$2^{O((\log n)^2)}$, and its exceptional boundary satisfies
$b=O((\log n)^2)$. Suppose all other charged work has the same
quasi-polynomial bound. With fresh independent sampling per call and
$\lambda=\lceil\log_2L(n)+c\log_2 n\rceil$, the total charged local
work is $n^{O(\log n)}$, and the probability that any required positive
call is missed is at most $n^{-c}$.
\end{corollary}
\begin{proof}
Substitute the hypotheses into Theorem~\ref{thm:main} and multiply by
$L(n)$. Both $2^{O(b)}$ and every fixed polynomial of the stated input size
are $2^{O((\log n)^2)}$. The chosen $\lambda$ is $O((\log n)^2)$ and
adds only a polynomial factor in the exponent's scale. Conditional on the
history before each call, that input is fixed and its fresh random variables
have miss probability at most $2^{-\lambda}$. The union bound is therefore
valid even for adaptively chosen later inputs. It gives
$L(n)2^{-\lambda}\leq n^{-c}$.
\end{proof}
The predecessor's $2^{O(b\log(b+1))}$ bound places the same local target,
using only that estimate, in the smaller sufficient regime
$b\log(b+1)=O((\log n)^2)$, for example
$b=O((\log n)^2/\log\log n)$. The new result removes this logarithmic
penalty in the parameter allowance, at the stated randomization cost.

The corollary is \emph{not} a complete unknot-recognition theorem. In
particular, the required feasible-forest hypothesis cannot be replaced by
``the local program usually succeeds.'' Some actual handlebody presentations
can stall under restricted cancellation rules \cite{repo,loop}, and some
abstract inputs have no successful forest at all. A general \Unk\ outcome
on such an input has probability one, not a small failure probability.
A correct integrated recognizer must continue through another verified route.

\begin{proposition}[Exact fallback accounting on positive local inputs]
Suppose a deterministic complete local fallback costs at most $F(\Lambda,b)$,
and $\lambda$ uncapped random rounds cost at most $\lambda K(\Lambda,b)$.
Running those rounds first and invoking the fallback on \Unk\ is always
correct. On a positive local input its expected work is at most
\[
 \lambda K(\Lambda,b)+2^{-\lambda}F(\Lambda,b).
\]
No corresponding small fallback probability follows for negative inputs.
\end{proposition}
\begin{proof}
Every early positive is independently valid. Otherwise the deterministic
fallback decides the local problem. On a positive input the probability of
reaching it is bounded by Theorem~\ref{thm:main}; on a negative input the
random positive search cannot succeed, so it may always be reached.
\end{proof}
Similarly, repeating random rounds until success has expected fewer than
two rounds on a promised positive instance, but need not terminate on a
negative one. Such a promise procedure is not a total recognition algorithm.
The finite-round public interface avoids that ambiguity.

The maintained supplied-cut representation has explicit combinatorial size
bounds in terms of tetrahedra and ports \cite{kernel}. Translating those
bounds into a polylogarithmic exceptional boundary across all useful
hierarchy stages remains a geometric theorem to be proved. Source surface
discovery, admissible regularization, passive port transport, pattern testing,
and restart charging still require their own bounds. Neither binary
compression nor a small local positive certificate proves those bounds.

\section{Further research questions and concrete next targets}\label{sec:future}
\subsection{Deterministic single-exponential mandatory contraction}
Can the sampled contraction and rank reduction be replaced by a deterministic
construction with $2^{O(b)}\poly(\Lambda)$ cost? A useful result must retain
mandatory extendibility while reducing rank; choosing arbitrary mandatory
representatives first can destroy a useful extra seed set. Fixing one
pseudorandom seed is not a solution. An explicit hitting family tailored to
the relevant graph minors, or a source-geometry restriction on their
interaction, would be a genuine strengthening of Theorem~\ref{thm:main}.
The unmixing proposition removes randomness from recovery, not from the
existence-preserving search representation.

\subsection{Source restrictions stronger than chain parity}
Characterize saturated matrix/dual-graph pairs realizable by the maintained
source charts. The parity-lift criterion is exact for its algebraic question
but only necessary for geometry. Can intersection ordering, vertex-link
planarity, bounded local source types, or boundary-port restrictions exclude
the abstract hard instances used in previous work? A construction of a
geometric realization of the nonlinear choice family would also clarify
which of its phenomena are genuinely topological. Conversely, a rigorous
nonrealizability proof would guide more meaningful benchmarks.

\subsection{A uniform geometric bound on the activation boundary}
The most direct route from the local theorem to the target complexity is a
bound $b=O((\log n)^2)$ for all required positive stages, together with the
other hypotheses of Corollary~\ref{cor:qp}. It may be more flexible to bound
the number of attainable closed masks rather than all $2^b$ masks, and to
exploit the smaller quotient rank $d$. Any such bound must survive the actual
hierarchy moves and not merely hold for one supplied cut.

\subsection{Joint optimization of primal and dual forests}
The one-handle basis is frozen throughout this article. Choosing a different
zero/one forest can change both column supports and the exceptional boundary.
Can a coupled selector trade primal changes against a smaller $b$ while
retaining independently replayable port charts? The graph-matroid part of
the current proof does not automatically encode these basis changes.
An end-to-end theorem would have to charge their effect on source occurrence
words and on all later matrix construction.

\subsection{More exact negative routes without full symbolic expansion}
The frontier certificate is complete when the representation is known to be
faithful. Can its coverage idea be lifted through a generic mandatory
contraction without certifying every symbolic minor? Possible directions
include rank-deficiency witnesses that explain all discarded mandatory
choices at once, or a deterministic source-specific representation.
A failed sampled frontier cannot be promoted to an original-input negative
certificate by confidence rhetoric.

\subsection{Fewer exceptional bits and stronger safe preprocessing}
Rows used only as heavy prerequisites can dominate $b$. An alternative
parameter based only on possible unit targets would need a compressed
representation of their prerequisite states; simply dropping the heavy rows
is invalid. Forced initial peeling and closure simplification may reduce the
relevant boundary, but the original dual edges must remain available as
forest resources even when their rows are already automatically active.
Deleting them from the dual graph without justification can change the
selection problem. A precise reduced interface for such preprocessing would
be useful both theoretically and in implementation.

\subsection{Weighted objectives and faster exterior computation}
For integration, not every successful tree is equally good: subsequent port
transport, word lengths, or proof size may matter. Weighted representative
families preserve certain additive costs, but a genuine downstream objective
need not be additive or determined by the current activation state. Identify
an objective with a valid extension-compatible invariant before pruning by
cost. Independently, implement sparse reachability, faster coordinate
insertion, and low-rank specializations; the present general dictionary code
is a research baseline rather than a tuned replacement for native sparse
primitives.

\subsection{Native replay, proof formalization, and adaptive sampling}
The immediate integration experiment is to run the proposal on authenticated
native source records, replay the source two/three certificate, and rerun
all downstream checks without transporting a stale local basis. The next
formalization target is the closure reduction and representative-extension
lemma, followed by the small deterministic frontier checker. Its lack of
dependence on producer algebra should make its trusted kernel manageable.
If samples are reused across adaptively generated local inputs, the simple
fixed-witness probability proof no longer applies automatically; fresh samples
or a stronger simultaneous guarantee must be provided.

\section{Conclusion}
The main result replaces a factorial/Bell-number search description by a
closed-state, rank-bounded representative search. Its decisive invariant is
that a successful dual tree contains at most $b$ useful extra seeds beyond
mandatory low-column representatives. Generic contraction handles the many
mandatory choices, and exterior span compression preserves every relevant
future independent extension. A fixed-witness probability argument avoids
an unnecessary simultaneous-faithfulness demand.

This yields a one-sided single-exponential local search, a deterministic
single-exponential low-rank subcase with checkable negative frontiers, and
explicit source-parity diagnostics. The mathematics, executable kernels,
independent checks, benchmarks, and integration boundaries are separated so
that the contribution can be incorporated without treating an inconclusive
local search as a knot verdict. The remaining global work is specific:
source-faithful geometry, bounded activation complexity over the hierarchy,
and complete charging of every required operator and restart.

\appendix
\section{The classical mandatory-representative oracle}\label{app:mi}
The prepass and current recovery implementation use graphic/partition
matroid intersection, specialized to disjoint color classes $U_i$. A fixed
forest $X$ is contracted logically: the independence oracle tests whether
$X\cup I$ is a forest. The partition constraint allows at most one selected
edge from each mandatory class.

Starting with $I=\varnothing$, form the exchange digraph on selectable
edges. An unselected edge is a source if it can be added to the forest
$X\cup I$. An unselected edge is a sink if its color is unused. A selected
edge $y$ points to an unselected $e$ when $X\cup(I-\{y\})\cup\{e\}$
is a forest. An unselected nonsink points to the selected edge of its color.
Toggle a shortest source-to-sink path. This is the classical common-
independence augmentation; in the graphic specialization its exchanges can
also be read directly from fundamental forest paths. The shortest-path
condition prevents earlier exchange chords from invalidating the toggle.
Iterating increases the common independent-set size until it is maximum.
The standard matroid-intersection and Rado argument is documented in the
audited predecessor \cite{prior}; this primitive is not asserted to be new.

For the uncontracted mandatory prepass, a failed augmentation yields a
subcollection of colors satisfying \eqref{eq:rado}. The implementation
returns that subcollection, and the independent checker recomputes the
mandatory components and graphic rank before accepting it. A maximum
matching count alone would not be an accepted certificate. For recovery
with fixed $X$, only actual representative edges are needed; their union
with $X$ and the resulting final tree are independently checked. The
16,038-query brute-force audit includes nonempty forced forests and empty
or deficient color classes.

\section{Source audit and contribution ledger}\label{app:audit}
The repository inspection is pinned, rather than silently following a
changing main branch. The maintained files inspected include the directory
README, the synthesis README, \path{synthesis/parallelity.tex},
\path{synthesis/bundle_boundary_kernel.tex},
\path{synthesis/core_three_handles.tex},
\path{synthesis/three_handles_results.tex}, and the complete current native
three-handle replay function. The corresponding source paths and available
Git blob identifiers are recorded in \path{SOURCE_AUDIT.json}.

The predecessor archive was read from the user's accessible Library and
materialized for programmatic comparisons. Its article and code were used
as a baseline; the vendored package retains its original license.
\path{ProveIt_Nonlinear_Activation_Forests_2026-10-12.zip} and
\path{ProveIt_Exceptional_Rows_2026-10-12.zip} were listed among the new
incoming files but their contents were not obtained. The ordered-interface
archive was inventoried, not used as a competing selector. Thus this study
claims a demonstrated advance relative to the inspected matroidal article,
not an exhaustive priority comparison against every incoming archive.

The ownership of the mathematical steps is also explicit. Private-column
completeness, the dual-forest source charts, the pair-column interpretation,
and the previous order/partition method come from the maintained and prior
materials. Generic linear specialization, finite-field polynomial testing,
graphic/partition intersection, and exterior representative families are
standard tools. The proved contribution here is their closed-activation
reduction and single-exponential combination, the faithful frontier
certificate protocol, the application-specific parity audit, and the stated
local quasi-polynomial consequences. The code, finite experiments, and
limits on their interpretation are part of that contribution. Neither
historical priority nor a completed global complexity theorem is asserted.

\begin{thebibliography}{99}
\bibitem{repo}
V.~Reshetnikov and ProveIt contributors.
\emph{Unknot Recognition: maintained project documentation}.
\path{Topology/UnknotRecognition/README.md}, pinned commit \pin, inspected
12 October 2026.
\url{https://github.com/VladimirReshetnikov/ProveIt}.

\bibitem{peeling}
ProveIt contributors.
\emph{Complete rectangular cancellation by sparse column peeling}.
\path{Topology/UnknotRecognition/synthesis/column_cancellation.tex}.
The result is also developed in the inspected preceding research article
\cite{prior} and referred to in the maintained kernel documentation.

\bibitem{dual}
ProveIt contributors.
\emph{Removing source core three-handles relative to the bundle ports}.
\path{Topology/UnknotRecognition/synthesis/core_three_handles.tex},
pinned commit \pin.

\bibitem{loop}
ProveIt contributors.
\emph{Source tests, dual-forest dependence, and three-handle results}.
\path{Topology/UnknotRecognition/synthesis/three_handles_results.tex},
pinned commit \pin.

\bibitem{kernel}
ProveIt contributors.
\emph{Parallelity geometry and the bundle-boundary cancellation kernel}.
\path{Topology/UnknotRecognition/synthesis/parallelity.tex} and
\path{Topology/UnknotRecognition/synthesis/bundle_boundary_kernel.tex},
pinned commit \pin.

\bibitem{prior}
ProveIt AI-assisted research contribution.
\emph{Matroidal Dual-Forest Selection for Unknot Recognition:
Polynomial pair-column cancellation, exact deletion deficits, and a
fixed-parameter higher-support kernel}.
12 October 2026. Article and executable code in
\path{docs/incoming/ProveIt_Matroidal_Dual_Forests_2026-10-12.zip}.
Its maintained-source checkpoint is
\texttt{c6d0ec8aa0b6308310ab6c081ba29d8bc363ede5}.

\bibitem{announcement}
Mathematical Institute, University of Oxford.
\emph{Marc Lackenby announces a new unknot recognition algorithm that runs
in quasi-polynomial time}. 2021; inspected 12 October 2026.
\url{https://www.maths.ox.ac.uk/node/38304}.

\bibitem{lackenby}
M.~Lackenby.
\emph{Incompressible surfaces, hierarchies and unknot recognition}.
arXiv:2607.23350v1, 25 July 2026, especially Sections 9--10.
\url{https://arxiv.org/abs/2607.23350}.

\bibitem{fomin}
F.~V.~Fomin, D.~Lokshtanov, F.~Panolan, and S.~Saurabh.
\emph{Efficient Computation of Representative Sets with Applications in
Parameterized and Exact Algorithms}.
arXiv:1304.4626v4, 2016.
\url{https://arxiv.org/abs/1304.4626}.

\bibitem{bevern}
R.~van Bevern, O.~Yu.~Tsidulko, and P.~Zschoche.
\emph{Representative families for matroid intersections, with applications
to location, packing, and covering problems}.
Discrete Applied Mathematics \textbf{298} (2021), 110--128.
arXiv:1806.11527v4.
\url{https://arxiv.org/abs/1806.11527}.

\bibitem{schwartz}
J.~T.~Schwartz.
\emph{Fast probabilistic algorithms for verification of polynomial identities}.
Journal of the ACM \textbf{27}(4) (1980), 701--717.
DOI: \href{https://doi.org/10.1145/322217.322225}{10.1145/322217.322225}.
\end{thebibliography}
\end{document}
